\documentclass{standalone}
\usepackage{pgfplots}
\pgfplotsset{compat=1.9}
\makeatletter
\newcommand{\clipprobe}{%
  \typeout{FINAL LIMITS: \pgfkeysvalueof{/pgfplots/xmin},\pgfkeysvalueof{/pgfplots/xmax},\pgfkeysvalueof{/pgfplots/ymin},\pgfkeysvalueof{/pgfplots/ymax}}%
  \pgfplotspointdescriptionxy{0}{0}\typeout{CLIP LOWER: \the\pgf@x,\the\pgf@y}%
  \pgfplotspointdescriptionxy{1}{1}\typeout{CLIP UPPER: \the\pgf@x,\the\pgf@y}%
  \pgfplotspointaxisxy{0}{0}\typeout{ORIGIN: \the\pgf@x,\the\pgf@y}%
  \pgfplotspointaxisxy{pi}{1}\typeout{DATA UPPER: \the\pgf@x,\the\pgf@y}%
}
\makeatother
\begin{document}
\begin{tikzpicture}
\begin{axis}[axis lines=middle,enlarge x limits=true,enlarge y limits=true,
  xmin=0,xmax=pi,ymin=-1,ymax=1,width=12cm,height=8cm,
  execute at end axis={\pgfplotsextra{\clipprobe}}]
\addplot[domain=-8:8,samples=200,blue,very thick] {cos(deg(x))};
\addplot[domain=-8:8,samples=200,red,very thick] {x};
\end{axis}
\end{tikzpicture}
\end{document}
